%!TEX root = ../main.tex

\chapter{Ist-Analyse und Anforderungen}

\section{Systemarchitektur von ATLAS aus Entwickler-Sicht}
Die Systemarchitektur umfasst Frontend, Backend, Datenhaltung und externe Dienste. Ziel dieses Abschnitts ist eine klare Darstellung der Komponenten und ihrer Schnittstellen.

\textbf{Zu beschreiben:}
\begin{itemize}
	\item Angular-Frontend, ASP.NET-Core-Backend, PostgreSQL und Azure Blob Storage
	\item Kommunikationswege zwischen Frontend und Backend
	\item Zweck zentraler Backend-Services
\end{itemize}

\section{Bestehende Testinfrastruktur}
Hier wird der Ausgangszustand der Tests beschrieben, einschließlich Umfang, Abdeckung und technischer Einbindung.

Im ersten Konzeptstand zeigte sich im Backend eine nur teilweise Abdeckung: Von fünf zentralen Services waren zunächst nur Experiment-, File- und Search-Service mit Tests versehen, und die Abdeckung lag insgesamt deutlich unter dem Zielwert von durchschnittlich 80\,\% in Kernbereichen. Für das Frontend lagen zu diesem Zeitpunkt noch keine automatisierten Tests und keine belastbare Coverage-Basis vor.

Zusätzlich war eine CI/CD-Ausführung zwar geplant (insbesondere als Pull-Request-nahe Qualitätssicherung), jedoch im betrachteten Zeitpunkt noch nicht aktiv.

\textbf{Zu beschreiben:}
\begin{itemize}
	\item bisheriger Test-Coverage-Anteil
	\item vorhandene Testarten und Werkzeuge
	\item Grenzen des aktuellen Ansatzes
\end{itemize}

\section{Anforderungen an die Testinfrastruktur (lokal und global)}
Die Anforderungen bilden die Grundlage für die Konzeption und spätere Bewertung der Lösung.

\subsection{Unittests}
Anforderung: schnelle, isolierte und lokal ausführbare Tests für zentrale Logikbausteine.

\subsection{Integrationstests}
Anforderung: Zusammenspiel mehrerer Komponenten mit realitätsnaher Daten- und Servicekopplung prüfen.

\subsection{End-to-End-Tests}
Anforderung: geschäftskritische Nutzerabläufe automatisiert über Systemgrenzen hinweg absichern.

\subsection{Klick- und UI-Tests}
Anforderung: robuste Validierung zentraler UI-Interaktionen und Frontend-Regressionen.

\subsection{Testumgebung}
Anforderung: reproduzierbare und wartbare Umgebungen für lokale und zentralisierte Ausführung.

\subsubsection{VM}
Option: definierte Laufzeitumgebung mit kontrollierten Abhängigkeiten.

\subsubsection{Azure DevOps}
Option: zentrale Pipeline-Orchestrierung für Build, Test und Reporting.

\subsubsection{Agenten}
Option: skalierbare Ausführung über dedizierte Build- oder Testagenten.

\subsubsection{Eigenständiger PC}
Option: dedizierte Hardware als pragmatische Zwischenlösung.

\subsection{Sicheres und schnelles Entwickeln}
Anforderung: kurze Feedbackzyklen bei gleichzeitig stabilen Qualitätskontrollen.

\subsection{Clean Code}
Anforderung: testspezifischer Code soll nachvollziehbar, konsistent und erweiterbar sein.

\subsection{Wartbarkeit}
Anforderung: langfristig pflegbare Teststruktur mit klaren Zuständigkeiten und geringem Anpassungsaufwand.

\subsection{Zusammenarbeit mit dem Hauptentwickler und bestehenden Tests}
Anforderung: neue Tests müssen sich konfliktarm in bestehende Testmuster und Entwicklungsprozesse integrieren.

\subsection{Funktionale und nichtfunktionale Anforderungen}
In diesem Unterabschnitt kann eine explizite Trennung in funktionale und nichtfunktionale Anforderungen vorgenommen und begründet werden.
